\begin{blocksection}
\question Complete \lstinline{make_fib} so the returned function remembers
Fibonacci results across calls. Assume non-negative integer inputs.
\begin{lstlisting}
def make_fib():
    cache = {0: 0, 1: 1}
    def fib(n):
        if __________________________:
            cache[n] = __________________________
        return __________________________
    return fib
\end{lstlisting}
\begin{solution}[0.5in]
\begin{lstlisting}
def make_fib():
    cache = {0: 0, 1: 1}
    def fib(n):
        if n not in cache:
            cache[n] = fib(n - 1) + fib(n - 2)
        return cache[n]
    return fib
\end{lstlisting}
\end{solution}
\begin{parts}
\part Execute these calls in order. Give each return value and the new
keys added to the cache during each call.
\begin{lstlisting}
f = make_fib()
f(5)
f(5)
f(6)
\end{lstlisting}
\begin{solution}[0.8in]
The first call returns 5 and adds keys 2, 3, 4, 5 in that order.
The second returns 5 and adds no keys: it is a cache hit.
The third returns 8 and adds only key 6; values for 5 and 4 are already stored.
\end{solution}
\part Why is the dictionary created outside \lstinline{fib}? Would
\lstinline{g = make_fib()} share \lstinline{f}'s saved results?
\begin{solution}[0.6in]
The returned function retains the same dictionary across recursive and
later calls. Creating it inside \lstinline{fib} would reset it on each
call. Calling \lstinline{make_fib} again creates a separate dictionary,
so \lstinline{g} does not share \lstinline{f}'s results.
\end{solution}
\end{parts}
\end{blocksection}
